\begin{center}\begin{tikzpicture}[font=\small,>=stealth,line join=round]
\draw[gray!40,rounded corners=3pt] (0,0) rectangle (14.2,8.3);\node at (7.1,7.700000000000001) {零函数同时具有两种对称};
\begin{scope}[shift={(4.1,3.9499999999999997)},x=1.25cm,y=2.3cm]\draw[->,gray] (-2.4,0)--(2.4,0) node[right]{$x$};\draw[->,gray] (0,-1)--(0,1) node[above]{$y$};\draw[blue,thick,solid,domain=-2:2,samples=180] plot(\x,{0});\fill[orange!85!black] (-1,0) circle(2.2pt);\fill[orange!85!black] (1,0) circle(2.2pt);\end{scope}\node[anchor=west,align=left,text width=5.0cm] at (8.4,4.0) {$H(x)=0$\\[10pt]$H(-x)=H(x)$\\[10pt]$H(-x)=-H(x)$\\[10pt]定义域仍须左右对应};
\node[text width=12.6cm,align=center] at (7.1,.65) {图像就是定义域对应的横轴部分；零高度既相等，也互为相反数。};
\end{tikzpicture}\end{center}
